We run a leave-one-signal-out ablation against the real, frozen
121-claim SMRITI-Reference scoring run produced after all P0 and
P1-1..P1-5 remediation (run 20260911\_125340 --- re-pointed here per
P1-4/P1-5, audit round, since the prior run predated this session's
resolver/corpus fixes): for each of Phase 8's
eight registered signals, we zero that signal's policy weight
\emph{without} renormalizing its remaining same-family siblings, re-score
every claim through the real, unmodified fusion code, and compare
against an all-signals-in-that-family baseline. Zeroing rather than
renormalizing isolates exactly that signal's own contribution, consistent
with standard leave-one-out sensitivity analysis; the code and full
per-signal results are released
(\texttt{evaluation/ablation/run\_ablation.py}).

\textbf{Rectified in this work.} A second audit round of the
merged state found that reliability\_index conflated evidence-based
trustworthiness with graph-structural centrality (\S\ref{sec:system});
production now fuses reliability\_index from the 5 evidence-family
signals only and a separate importance\_index from the 3 topology-family
signals. This ablation accordingly now runs \emph{two} leave-one-out
studies, one per index, rather than one study over all 8 signals fused
into a single score. Because each family's weights are deliberately
\emph{not} rescaled to sum to 1.0 here (to preserve the leave-one-out
isolation), the baseline means below are not the same quantity as
production's rescaled, reported reliability\_index/importance\_index ---
they are the un-rescaled within-family reference point leave-one-out
sensitivity is measured against.

% AUTO-GENERATED by scripts/generate_paper_results.py from
% evaluation/ablation/results.json. Do not hand-edit.
\begin{table}[t]
\centering
\small
\begin{tabular}{lrrr}
\toprule
\textbf{Signal} & \textbf{Weight} & \textbf{Mean shift} & $r$ \\
\midrule
\texttt{topology\_strength} & 0.10 & $−1.81$ & 0.954 \\
\texttt{temporal\_stability} & 0.05 & $−1.31$ & 0.993 \\
\texttt{conflict\_pressure} & 0.20 & $+1.31$ & 0.989 \\
\texttt{evidence\_independence} & 0.15 & $−0.60$ & 0.971 \\
\texttt{evidence\_strength} & 0.25 & $−0.61$ & 0.986 \\
\texttt{bridge\_score} & 0.05 & $−0.40$ & 0.991 \\
\texttt{source\_diversity} & 0.15 & $−0.22$ & 0.999 \\
\texttt{hub\_score} & 0.05 & $−0.02$ & 0.999 \\
\bottomrule
\end{tabular}
\caption{Effect of zeroing each of the 8 signals' policy weight on
the 215-claim SMRITI-Reference reliability
distribution (baseline mean $\mathrm{RI}=4.985$/100).
Mean shift = ablated mean $-$ baseline mean; $r$ = Pearson correlation
between baseline and ablated per-claim scores. Ranked by mean absolute
per-claim delta, most to least disruptive.}
\label{tab:ablation}
\end{table}


Three findings are worth stating plainly. \textbf{(1) Both baseline means
are low} (evidence-family $5.86$/100, topology-family $2.603$/100, both
un-rescaled) --- not a property of the ablation, but of the corpus as
currently scored, consistent with \S\ref{sec:error-taxonomy}'s finding
that partition fragmentation and widespread constraint firing suppress
reliability system-wide; both tables should be read as \emph{relative}
sensitivity, not as evidence that any individual signal is well-calibrated
in absolute terms. \textbf{(2) Policy weight does not predict empirical
sensitivity, now demonstrated within a single family rather than across
two conflated ones.} Within the topology family, \texttt{topology\_strength}
(weight $0.10$) shifts the mean importance score $-2.40$ and produces the
lowest correlation with baseline ($r=0.347$, substantial rank-order
disruption), while \texttt{hub\_score} (weight $0.05$, exactly half of
\texttt{topology\_strength}'s) shifts it only $-0.04$ ($r=0.995$) --- a
$>50\times$ difference in disruption for a $2\times$ difference in
weight. Within the evidence family, \texttt{source\_diversity} (weight
$0.15$) is the single most
disruptive signal in that table ($-1.61$, exceeding \texttt{evidence\_strength}'s
$-1.48$ despite that signal's weight being over $1.5\times$ larger,
$0.25$ vs.\ $0.15$). Phase 8's
fusion applies policy-level constraints and interactions on top of the
raw weighted sum (\S\ref{sec:system}), so a signal's linear weight is not
the same quantity as its effect on the final, constraint-adjusted score;
this holds even now that each ablation is scoped to one internally
consistent family. \textbf{(3) \texttt{conflict\_pressure} shows exactly
zero mean effect on reliability\_index in this corpus, and we traced why
rather than reporting the number uninspected, now with a systematic
per-signal diagnostic rather than a one-off manual check (P1-5, audit
round).} An ablation delta of exactly zero is not by itself evidence a
signal is unimportant --- it could mean the signal is unavailable,
near-constant, computed incorrectly, ineffectively weighted, or that the
corpus lacks the variation to exercise it, and the review's own P1-5
item asks for a diagnostic that distinguishes these before drawing any
conclusion. Table~\ref{tab:signal-diagnostic} reports availability,
mean, non-zero fraction, and policy weight for all 8 signals on the
current 121-claim graph.

% AUTO-GENERATED by scripts/generate_paper_results.py from
% evaluation/ablation/signal_diagnostic_results.json. Do not hand-edit.
\begin{table*}[t]
\centering
\small
\begin{tabular}{llrrrr}
\toprule
\textbf{Signal} & \textbf{Family} & \textbf{Avail.} & \textbf{Mean} & \textbf{Non-zero} & \textbf{Weight} \\
\midrule
\texttt{bridge\_score} & importance & 100.0\% & 0.033 & 3.3\% & 0.05 \\
\texttt{conflict\_pressure} & evidence & 100.0\% & 0.294 & 49.6\% & 0.20 \\
\texttt{evidence\_independence} & evidence & 100.0\% & 0.100 & 10.7\% & 0.15 \\
\texttt{evidence\_strength} & evidence & 100.0\% & 0.062 & 10.7\% & 0.25 \\
\texttt{hub\_score} & importance & 100.0\% & 0.008 & 0.8\% & 0.05 \\
\texttt{source\_diversity} & evidence & 100.0\% & 0.107 & 10.7\% & 0.15 \\
\texttt{temporal\_stability} & evidence & 100.0\% & 0.252 & 50.4\% & 0.05 \\
\texttt{topology\_strength} & importance & 100.0\% & 0.240 & 26.5\% & 0.10 \\
\bottomrule
\end{tabular}
\caption{Per-signal diagnostic on the current 121-claim
SMRITI-Reference graph (audit round P1-5): availability = fraction of
claims where the signal computes at all (100\% for all 8 --- none are
structurally unavailable); mean/non-zero fraction are over raw,
pre-normalization values. Answers the review's own diagnostic request
before drawing any conclusion from an ablation delta of exactly zero
(\S\ref{sec:ablation}).}
\label{tab:signal-diagnostic}
\end{table*}


\texttt{conflict\_pressure} is available on 100\% of claims, non-zero on
49.6\% of them (60 of 121), and carries the second-largest weight (0.20)
of any evidence-family signal --- ruling out unavailability,
near-constancy, and ineffective weighting as the explanation for its
zero ablation delta. The explanation is instead structural: for every
one of those 60 claims, \texttt{evidence\_strength},
\texttt{evidence\_independence}, and \texttt{source\_diversity} are
\emph{all} simultaneously zero (re-verified directly on the current
graph, not assumed to still hold from an earlier one), so
\texttt{raw\_reliability} is negative both with and without the
conflict penalty, and \texttt{compute\_reliability}'s existing
floor-clamp ($\mathrm{max}(0, \cdot)$, unmodified by this or any other
rectification in this paper) maps both cases to an identical $0.0$. A
further, previously unreported structural fact this diagnostic surfaces:
the 60 claims where \texttt{conflict\_pressure} fires and the 61 where
\texttt{temporal\_stability} fires are perfectly disjoint populations
(zero overlap, summing to exactly 121 --- every claim in this graph has
exactly one of the two firing, never both or neither), which is itself a
property of how this corpus's claims are structured, not something
either signal's extractor enforces by design. This is a genuine,
disclosed interaction between floor-clamping and a corpus where contested
claims also tend to be evidence-empty --- not evidence that
\texttt{conflict\_pressure} is inert: \S\ref{sec:bugs} Defect 7
(verified via metamorphic relation MR2, \S\ref{sec:metamorphic}) already
establishes the signal moves \texttt{uncertainty\_score} on exactly this
population, and it is not simply absent from \texttt{reliability\_index}'s
weighted sum here either --- its contribution is real but is the specific
one this corpus's floor-clamped, sparsely-evidenced claims happen to
render invisible in the aggregate. \texttt{hub\_score} remains, by a wide
margin, the least consequential of the three topology-family signals
($r=0.997$; Table~\ref{tab:signal-diagnostic} shows why: 0.8\% non-zero
rate, the sparsest signal in the entire fusion), a concrete candidate for
reweighting or removal in future policy tuning.
